\chapter{Der Übergang und seine Richtung}

\section{Zwei entgegengesetzte Übergänge}

Zwei der stärksten Gedanken, die von Zeit und Raum handeln, behaupten Übergänge in entgegengesetzter Richtung. Hegel: Der Raum geht in die Zeit über. Die Wahrheit des Raumes ist die Zeit; der Punkt, der im Raum keine Wirklichkeit hat, hat sie in der Zeit; \zitat{die Zeit ist also erst der wahrhafte Punkt}. Guzzoni, mit Heidegger, Wagner und Hölderlin: Die Zeit wird zum Raum. Sie eilt hin zum Ort, sie weilt und gewinnt Fläche, sie ist immer schon eingeräumt. Wer beide Sätze nebeneinanderstellt, sieht einen Widerspruch. Wer fragt, wovon sie jeweils sprechen, sieht keinen.

Hegels Übergang ist ein begrifflicher. Das Abstraktere hat seine Wahrheit im Konkreteren. Der Raum, das gleichgültige Nebeneinander, ist abstrakter als die Zeit, das Werden; deshalb ist die Zeit seine Wahrheit. Der Übergang sagt nichts darüber, was in der Erfahrung zuerst da ist oder woraus etwas entsteht; er ist kein Vorgang. Bonsiepen hat gezeigt, dass Hegel seit der Dissertation von 1801 dem abstrakten Raum die Subjektivität hinzufügen muss, den Punkt, die Zeit, und dass der Punkt, im Jenaer Systementwurf, \zitat{mens} heißt: Der Übergang vom Raum in die Zeit ist der Übergang von der Sache zum Geist.\footnote{Bonsiepen, Hegels Raum-Zeit-Lehre, zur Genese.}

Guzzonis Übergang ist einer der Erfahrung. Wenn die Zeit weilt, wird sie Raum: Sie breitet sich aus, sie hat Gegend, sie lässt sich aufhalten. Der Ort ist der Augenblick, an dem die Zeit sich niedergelassen hat. Dieser Übergang ist ein Vorgang, aber kein zeitlicher Vorgang in der Zeit; er ist die Weise, wie die Zeit ist, wenn sie nicht rennt.

\begin{figure}[htbp]
\centering
\begin{tikzpicture}[scale=1,>=Stealth,every node/.style={font=\small}]
% obere Reihe: Hegels räumlicher Aufstieg
\node (p) at (0,2) {Punkt};
\node (l) at (2.6,2) {Linie};
\node (f) at (5.2,2) {Fläche};
\node (r) at (7.8,2) {Raum};
\draw[->] (p) -- (l); \draw[->] (l) -- (f); \draw[->] (f) -- (r);
\node[gray] at (3.9,2.55) {\scriptsize der Punkt hat im Raum keine Wirklichkeit};
% Übergang Raum -> Zeit
\node (z) at (7.8,0) {Zeit};
\draw[->,thick] (r) -- node[right] {\scriptsize Wahrheit des Raumes} (z);
% untere Reihe: Zeit -> Jetzt -> Ort
\node (j) at (5.2,0) {Jetzt};
\node (o) at (2.6,0) {Ort};
\node (g) at (0,0) {Gegend};
\draw[->] (z) -- node[above] {\scriptsize \zitat{der wahrhafte Punkt}} (j);
\draw[->] (j) -- node[above] {\scriptsize räumliches Jetzt} (o);
\draw[->] (o) -- node[above] {\scriptsize Weile} (g);
% Rückkehr: Ort ist im Raum
\draw[->,thick,dashed] (o) -- node[left] {\scriptsize \zitat{immer schon eingeräumt}} (p);
% Beschriftung der Richtungen
\node[gray,align=left] at (10.4,2) {\scriptsize Hegel:\\[-2pt]\scriptsize der Raum geht\\[-2pt]\scriptsize in die Zeit über};
\node[gray,align=left] at (10.4,0) {\scriptsize Guzzoni:\\[-2pt]\scriptsize die Zeit wird\\[-2pt]\scriptsize zum Raum};
\end{tikzpicture}
\caption{Der Übergang und seine Richtung. Oben der begriffliche Aufstieg, in dem der Punkt sich im Raum ausbreitet und der Raum in die Zeit übergeht (Hegel). Unten der Weg der Erfahrung, in dem die Zeit als Jetzt zum Ort und als Weile zur Gegend wird (Guzzoni). Die gestrichelte Rückkehr: Der Ort ist ein Punkt im Raum, die Zeit ist eingeräumt.}
\label{abb:uebergang}
\end{figure}

\section{Dieselbe Rangordnung}

Beide Übergänge sagen dasselbe über den Rang. Bei Hegel ist die Zeit das Wahrhafte, der Raum die \zitat{leere Zerstreuung des Außersichseins}. Bei Guzzoni ist die Zeit \zitat{das Vorrangige und Bestimmende}, der Raum das, worin sie eingeräumt ist. Bei Reichenbach ist die Zeit das logisch Primäre, der Raum das Nebeneinander der Fasern. Bei Brentano ist der Begriff des Zeitlichen dem des Räumlichen übergeordnet. Bei Chiba ist die Zeit die weitere Klasse. Bei Baumann ist die Dauer des Ich das, ohne welches die Zeitvorstellung nicht gebildet werden könnte, und der Raum eine Anschauung, die uns ohne Dinge bleibt, aber nichts hervorbringt.

Die Übereinstimmung geht weiter als der Rang. Sie betrifft die \emph{Richtung der Ableitbarkeit}. Aus der Zeit erhält man den Raum: Die Gleichzeitigkeit gibt das Nebeneinander, das Jetzt gibt den Ort, die Lichtlaufzeit gibt die Länge, das Nahwirkungsprinzip gibt die Dimensionszahl. Aus dem Raum erhält man keine Zeit. Zenon scheitert daran, den Übergang aus Orten zusammenzusetzen. Brentanos Zeit mit einem einzigen Präteritalmodus wäre ein Topoid, also Raum. Und Leibniz' Definition des Raumes als Ordnung des Zugleichseienden ist, wie Baumann gezeigt hat, zirkulär: Das Zugleich ist zeitlich, und um daraus den Raum zu gewinnen, muss man \zitat{den Raum bereits mitbringen}; \zitat{Im Grunde lautet die Definition: Raum ist die Ordnung des Räumlichen.}\footnote{Baumann, Lehren, Bd.\,2, Kapitel zu Leibniz' Raum- und Zeitlehre.} Man kann also den Raum durch die Zeit definieren, und die Definition ist leer, weil sie das Räumliche voraussetzt; und man kann die Zeit nicht durch den Raum definieren, weil die Definition ein Topoid ergibt. Das ist die Asymmetrie in ihrer logischen Form.

\section{Der Zeitraum}

Kant hat in der Analytik beides gesehen. Martić fasst zusammen: Der Raum setzt die Zeit voraus und ist nur als Zeitraum konzipierbar; die Zeit wird nur am Raum zur Größe. Er zitiert Schelling: \zitat{Die Zeit wird nur durch den Raum, der Raum nur durch die Zeit endlich.}\footnote{Martić, Zeit und Raum, Kapitel zum Schematismus und zur Rezeption bei Fichte und Schelling.} Das Wechselverhältnis ist also nicht Parallelisierung. Es ist eine Ordnung mit einer Richtung und einem Vorbehalt. Die Richtung: Die Zeit stiftet die Ordnung, der Raum ist Zeitraum. Der Vorbehalt: Ohne Raum keine Größe, ohne Beharrliches keine Bestimmung, ohne Messkörper keine Lichtgeometrie.

Dass die beiden Übergänge sich zu diesem Wechselverhältnis fügen, ist das Ergebnis. Hegels Aufstieg vom Raum zur Zeit sagt, was das Wahrhafte ist; Guzzonis Weg von der Zeit zum Ort sagt, wo es sich aufhält. Beide zusammen sagen: Die Zeit ist das, was der Raum in Wahrheit ist, und der Raum ist das, worin die Zeit sich aufhält. Keine der beiden Seiten kann ohne die andere sein. Aber sie sind nicht gleich. Die eine gibt, die andere nimmt auf.
